\section{THE THEORY OF LIE ALGEBRAS.}

Once the correspondence between groups of transformations and Lie algebras is acquired, the theory is going to take a clearly more algebraic turn and will be centred on a deep study of Lie algebras. \(^{13}\)

A first and short period, from 1888 to 1894, marked by the work of Engel, of his pupil Umlauf and especially Killing and E. Cartan, ended up with a series of spectacular results on complex Lie algebras. We have seen above that the notion of soluble Lie algebra was due to Lie himself, who had proved (in the complex case) the reduction theorem for soluble linear Lie algebras to triangular form ({[}203{]}, v. I, p.~270). \(^{14}\) Killing observes {[}180{]} that there exists in a Lie algebra a biggest soluble ideal (that is called the radical today), and that the quotient of the Lie algebra by its radical has null radical; he calls the Lie algebras with null radical semisimple, and proves that they are the product of simple algebras (this last notion had already been introduced by

Lie, who had proved the simplicity of the ``classical'' Lie algebras ({[}203{]}, v. 3, p.~682)).

On the other hand, Killing introduces, in a Lie algebra, the characteristic equation \(\det(\operatorname{ad}(x)-\omega.1)=0\) , already met by Lie in studying the Lie subalgebras of dimension 2 containing a given element of a Lie algebra. We refer to other historical Notes for the analysis of the methods whereby Killing, in studying in a penetrating manner the properties of the roots of the ``generic'' characteristic equation for a semisimple algebra, ends up with the most remarkable of his results, the complete determination of the simple (complex) Lie algebras. \(^{15}\)

Killing proves that the derived algebra of a soluble algebra is ``of rank 0'' (which means that ad \(x\) is nilpotent for every element \(x\) of the algebra). Soon after, Engel proves that the algebras ``of rank 0'' are soluble (this statement is in substance what is called today Engel's theorem). In his thesis, E. Cartan introduces on the other hand what is now called the ``Killing form'', and establishes the two fundamental criteria that characterise by means of this form the soluble Lie algebras and the semisimple Lie algebras.

Killing had affirmed ({[}180{]}, IV) that the derived algebra of a Lie algebra is the sum of a semisimple algebra and of its radical, which is nilpotent, but his proof was incomplete. A little later, E. Cartan announced without proof ({[}52 a{]}, v. I₁, p.~104) that more generally every Lie algebra is the sum of its radical and of a semisimple subalgebra; the only result in this direction that is indisputably established at this time is a theorem of Engel affirming the existence, in every non-soluble Lie algebra, of a simple Lie subalgebra of dimension 3. The first published proof (for complex Lie algebras) of the statement of Cartan is due to E. E. Levi {[}200{]}; another proof (equally valid for the real case) was given by J. H. C. Whitehead in 1936 {[}334 a{]}. In 1942, A. Malcev completed this result by the uniqueness theorem on the ``Levi sections'' up to conjugation.

Already from his first work, Lie had set himself the problem of the isomorphism of every Lie algebra with a linear Lie algebra. He had thought that he had solved it affirmatively by considering the adjoint representation (and so deducing a proof of his ``third theorem'') ({[}202{]}, vol.~5, Abh. III, pp.~42-75); he recognised rapidly that his proof was only correct for Lie algebras with null centre; after him, the question remained open for a long time, and was solved affirmatively by Ado in 1935 {[}2 a{]}. On the other hand, Lie had substantially set himself the problem of determining the linear representations of minimal dimension of the simple Lie algebras, and had solved it for the classical algebras; in his Thesis, Cartan had also resolved this problem for the exceptional simple algebras; \(^{16}\) the methods that he employs in this context will be generalised by him twenty years later in order to obtain all the irreducible representations of simple real or complex Lie algebras.

The property of complete reducibility of a linear representation seems to have been met for the first time (in a geometric form) by Study. In an unpublished manuscript, but quoted in ({[}203{]}, v. 3, pp.~785-788) he proves this property for the linear representation of the Lie algebra of SL(2, C), and obtains partial results for SL(3, C) and SL(4, C). Lie and Engel conjecture on this occasion that the complete reducibility theorem is valid for SL(n, C) whatever n may be. The complete reducibility of the linear representations of semisimple Lie algebras was established by H. Weyl in 1925 \(^{17}\) by an argument of a global nature (see later). The first algebraic proof was obtained in 1935 by Casimir and van der Waerden {[}55{]}; other algebraic proofs were given afterwards by R. Brauer {[}34 b{]} and J. H. C. Whitehead {[}334 b{]}.

Finally, during his research on the exponential map (cf.~infra). H. Poincaré ({[}251{]}, v. III) considers the associative algebra of differential operators of all orders, generated by the operators of a Lie algebra; he shows, substantially that, if \((X_{i})_{1\leq i\leq n}\) is a basis of a Lie algebra, the associative algebra generated by the \(X_{i}\) has as basis certain symmetric functions of the \(X_{i}\) (sums of the non-commutative ``monomials'' obtained from a given monomial by all permutations of its factors). The essential part of his proof is algebraic in nature, and allows him to obtain the structure of the enveloping algebra. Analogous proofs have been given in 1937 by G. Birkhoff {[}22 b{]} and E. Witt {[}337 b{]}. \(^{18}\)

The majority of the work quoted above is limited to real or complex Lie algebras, which are the only ones corresponding to Lie groups in the usual sense. The study of Lie algebras over a field other than R or C is tackled by Jacobson {[}172 a{]} who shows that the greater part of the classical results remain valid over a field of characteristic zero.

